\section{The bilinear inequality and equality cases}
\label{sec:conclusion}

We now combine the spherical and variational estimates. Recall that an
admissible measure satisfies the support-functional inequality
\eqref{eq:functional-condition}, and need not itself be symmetric.

\begin{theorem}\label{thm:bilinear}
Let $\eta,\zeta$ be finite positive measures on $\R^n$ with bounded,
spanning support. Suppose that both satisfy
\eqref{eq:functional-condition}. Then
\begin{equation}\label{eq:bilinear}
 \iint |x\cdot y|\,d\eta(x)d\zeta(y)\ge b_n.
\end{equation}
If equality holds, there is an origin-centered ellipsoid $E_0$ such that
\[
 \int h_{E_0}\,d\eta=(|E_0|/\kappa_n)^{1/n}.
\]
\end{theorem}

\begin{proof}
Choose $k$ as in \eqref{eq:k-choice} and write $c=c_{n,k}$.
Apply Proposition~\ref{prop:position} to $\eta$, obtaining a symmetric
positive definite matrix $A$ with determinant one.
Replace $\eta$ by its image under $A$ and $\zeta$ by its image under
$A^{-1}$. Their pairing is unchanged, since
$(Ax)\cdot(A^{-1}y)=x\cdot y$.
Their support-functional inequalities also persist: for any linear map
$T$,
\[
 h_M(Tx)=h_{T^tM}(x),
\]
and $|T^tM|=|M|$ when $|\det T|=1$.

For the first transformed measure, the support functional at the
identity is precisely the original $\tau_A$. Its unique minimizing
body is therefore $L_A$, with the degree-two cancellation supplied by
Proposition~\ref{prop:position}.
For each transformed measure choose its minimizing body as in
Proposition~\ref{prop:representation}, and rescale that body to volume
$\kappa_n$. Let the resulting gauges be $g_1,g_2$, and put
\[
 f_i=g_i^k,\qquad t_i=\E g_i^{k-1},\qquad
 s_i=\frac{\tau_i(L_i)}{t_i},\qquad
 a_i=s_i\E f_i\quad(i=1,2),
\]
where $\tau_i$ is the functional of the corresponding transformed
measure. Thus
\begin{equation}\label{eq:normalized-parameters}
 \E g_i^{-n}=1,\qquad s_it_i\ge1,\qquad a_i\ge1.
\end{equation}
For the last inequality, if $k=1$ then $\E g_i\ge1=t_i$ by the
negative-moment normalization. If $k=2$, the same normalization gives
$\E g_i\ge1$, and $\E g_i^2\ge(\E g_i)^2\ge\E g_i=t_i$.
The first function $f_1$ has no degree-two harmonic component; scalar
rescaling does not alter this property.

Apply \eqref{eq:representation} to $|x\cdot y|$ successively in both
variables. These are even continuous homogeneous tests, so the formula
applies despite any asymmetry of the measures. It gives
\begin{align}
 \iint |x\cdot y|\,d\eta(x)d\zeta(y)
 &=s_1s_2\iint
       |X_{f_1}(u)\cdot X_{f_2}(v)|\,d\sigma(u)d\sigma(v)
       \notag\\
 &\ge b_ns_1s_2\mathcal B(f_1,f_2).
 \label{eq:pairing-sign}
\end{align}
The pointwise inequality in the last line uses the test
$\sgn(u\cdot v)$. Proposition~\ref{prop:sign-estimate} now gives
\[
 b_n^{-1}\iint |x\cdot y|\,d\eta(x)d\zeta(y)
 \ge a_1a_2-s_1s_2c\norm{Qf_1}_2\norm{Qf_2}_2.
\]
By Theorem~\ref{thm:norm-estimate} and
\eqref{eq:normalized-parameters},
\begin{equation}\label{eq:error-absorption}
 s_i^2c\norm{Qf_i}_2^2
 \le a_i^2-s_i^2t_i^2\le a_i^2-1.
\end{equation}
Consequently the last lower bound is at least
\[
 a_1a_2-\sqrt{(a_1^2-1)(a_2^2-1)}\ge1.
\]
Indeed $a_1a_2-1\ge0$ and
\[
 (a_1a_2-1)^2-(a_1^2-1)(a_2^2-1)=(a_1-a_2)^2\ge0.
\]
This proves \eqref{eq:bilinear}.

Suppose equality holds. If $a_i>1$, the inequality
$s_i^2c\norm{Qf_i}_2^2\le a_i^2-1$ is strict: this follows from
Theorem~\ref{thm:norm-estimate} when $g_i$ is nonconstant, and from its
zero left side when $g_i$ is constant. If both $a_i$ exceed one, the
product error in \eqref{eq:error-absorption} is therefore strictly
smaller than the displayed square root. If exactly one exceeds one,
the error product vanishes and $a_1a_2>1$.
Both alternatives contradict equality, so $a_1=a_2=1$.
In particular, $\E g_1^k\le\E g_1^{k-1}$.
For $k=1$, this forces $\E g_1=1$ and equality in the negative-moment
power-mean bound; hence $g_1=1$.
For $k=2$, the chain
\[
 1\le\E g_1\le(\E g_1)^2\le\E g_1^2\le\E g_1
\]
again forces constancy and then $g_1=1$.
Thus $s_1=1$, and the support-functional condition of the transformed
first measure is sharp on $B_2^n$. Undoing the transformation makes the
original condition sharp on $A^tB_2^n$, which is an origin-centered
ellipsoid of volume $\kappa_n$.
\end{proof}

\begin{proof}[Proof of Theorem~\ref{thm:main}]
Translate $K$ so that the origin is in its interior, and first rescale
it to satisfy $|K|=\kappa_n$.
Proposition~\ref{prop:projection-measure} supplies the admissible
measure $\eta_K$. Set
\[
 r=\left(\frac{|\Pi K|}{\kappa_n}\right)^{1/n},
 \qquad D=\frac1r\Pi K.
\]
The symmetric body $D$ has volume $\kappa_n$, so it too has an
admissible measure $\eta_D$. Its own support-functional integral is
one: the volume derivative under $D+tD=(1+t)D$ gives
\[
 \int h_D\,d\eta_D
 =\frac1{n\kappa_n}\left.\frac{d}{dt}\right|_{0+}|(1+t)D|=1.
\]
Using \eqref{eq:projection-measure} and Theorem~\ref{thm:bilinear},
\[
 1=\frac{n\kappa_n}{2r}
       \iint|x\cdot y|\,d\eta_K(x)d\eta_D(y)
   \ge\frac{n\kappa_nb_n}{2r}.
\]
Equation~\eqref{eq:ball-constant} yields $r\ge\kappa_{n-1}$, or
$|\Pi K|\ge\kappa_n\kappa_{n-1}^n$ under the current normalization.
The projection body scales as $\Pi(aK)=a^{n-1}\Pi K$ for $a>0$.
Undoing the volume normalization proves \eqref{eq:main}.

If equality holds, the pairing of $\eta_K$ and $\eta_D$ equals $b_n$.
Theorem~\ref{thm:bilinear} gives an ellipsoid on which the
support-functional bound for $\eta_K$ is sharp.
The equality assertion of Proposition~\ref{prop:projection-measure}
then makes $K$ a translate of a positive dilate of that ellipsoid.

For the converse, the quotient is invariant under translations and
invertible linear maps. To verify the latter directly, the segment
volume derivative gives, for every vector $y$ and every invertible $T$,
\begin{align*}
 2h_{\Pi(TK)}(y)
 &=\left.\frac{d}{dt}\right|_{0+}|TK+t[-y,y]|\\
 &=|\det T|\left.\frac{d}{dt}\right|_{0+}
              |K+t[-T^{-1}y,T^{-1}y]|\\
 &=2|\det T|h_{\Pi K}(T^{-1}y).
\end{align*}
Hence $\Pi(TK)=|\det T|T^{-t}\Pi K$ and
$|\Pi(TK)|=|\det T|^{n-1}|\Pi K|$, exactly matching the factor in
$|TK|^{n-1}$. Finally, $\Pi B_2^n=\kappa_{n-1}B_2^n$ gives equality
for the ball, and therefore for every ellipsoid.
\end{proof}
